\documentclass[10pt,a4paper]{amsart}
\usepackage{amsmath,amssymb,amsthm,graphicx}
% For commutative diagrams we use Xy-pic
\usepackage[all]{xy}

% We use 2cell for 2-commutative diagrams.
\xyoption{2cell}
% Theorem environments.
%
\theoremstyle{plain}
\newtheorem{theorem}[subsection]{Theorem}
\newtheorem{proposition}[subsection]{Proposition}
\newtheorem{lemma}[subsection]{Lemma}

\theoremstyle{definition}
\newtheorem{definition}[subsection]{Definition}
\newtheorem{example}[subsection]{Example}
\newtheorem{exercise}[subsection]{Exercise}
\newtheorem{situation}[subsection]{Situation}

\theoremstyle{remark}
\newtheorem{remark}[subsection]{Remark}
\newtheorem{remarks}[subsection]{Remarks}

\numberwithin{equation}{subsection}

% Macros
%
\def\lim{\mathop{\mathrm{lim}}\nolimits}
\def\colim{\mathop{\mathrm{colim}}\nolimits}
\def\Spec{\mathop{\mathrm{Spec}}}
\def\Hom{\mathop{\mathrm{Hom}}\nolimits}
\def\Ext{\mathop{\mathrm{Ext}}\nolimits}
\def\SheafHom{\mathop{\mathcal{H}\!\mathit{om}}\nolimits}
\def\SheafExt{\mathop{\mathcal{E}\!\mathit{xt}}\nolimits}
\def\Sch{\mathit{Sch}}
\def\Mor{\mathop{\mathrm{Mor}}\nolimits}
\def\Ob{\mathop{\mathrm{Ob}}\nolimits}
\def\Sh{\mathop{\mathit{Sh}}\nolimits}
\def\NL{\mathop{N\!L}\nolimits}
\def\CH{\mathop{\mathrm{CH}}\nolimits}
\def\proetale{{pro\text{-}\acute{e}tale}}
\def\etale{{\acute{e}tale}}
\def\QCoh{\mathit{QCoh}}
\def\Ker{\mathop{\mathrm{Ker}}}
\def\Im{\mathop{\mathrm{Im}}}
\def\Coker{\mathop{\mathrm{Coker}}}
\def\Coim{\mathop{\mathrm{Coim}}}

% Boxtimes
%
\DeclareMathSymbol{\boxtimes}{\mathbin}{AMSa}{"02}

%
% Macros for moduli stacks/spaces
%
\def\QCohstack{\mathcal{QC}\!\mathit{oh}}
\def\Cohstack{\mathcal{C}\!\mathit{oh}}
\def\Spacesstack{\mathcal{S}\!\mathit{paces}}
\def\Quotfunctor{\mathrm{Quot}}
\def\Hilbfunctor{\mathrm{Hilb}}
\def\Curvesstack{\mathcal{C}\!\mathit{urves}}
\def\Polarizedstack{\mathcal{P}\!\mathit{olarized}}
\def\Complexesstack{\mathcal{C}\!\mathit{omplexes}}
% \Pic is the operator that assigns to X its picard group, usage \Pic(X)
% \Picardstack_{X/B} denotes the Picard stack of X over B
% \Picardfunctor_{X/B} denotes the Picard functor of X over B
\def\Pic{\mathop{\mathrm{Pic}}\nolimits}
\def\Picardstack{\mathcal{P}\!\mathit{ic}}
\def\Picardfunctor{\mathrm{Pic}}
\def\Deformationcategory{\mathcal{D}\!\mathit{ef}}
\usepackage{cleveref}
\crefname{originallemma}{original lemma}{original lemmas}
\Crefname{originallemma}{Original lemma}{Original lemmas}
\makeatletter
\newlabel{lemma-locally-constant-on-connected}{{0DV5}{source}{}{sourcepointer-1}{}}
\newlabel{lemma-locally-constant-on-connected@cref}{{[originallemma][0][]0DV5}{[1][0][]source}{}{}{}}
\makeatother
\crefname{originallemma}{original lemma}{original lemmas}
\Crefname{originallemma}{Original lemma}{Original lemmas}
\makeatletter
\newlabel{morphisms-lemma-normalization-geom-unibranch-univ-homeo}{{0GIQ}{source}{}{sourcepointer-2}{}}
\newlabel{morphisms-lemma-normalization-geom-unibranch-univ-homeo@cref}{{[originallemma][0][]0GIQ}{[1][0][]source}{}{}{}}
\makeatother
\crefname{originallemma}{original lemma}{original lemmas}
\Crefname{originallemma}{Original lemma}{Original lemmas}
\makeatletter
\newlabel{pione-proposition-universal-homeomorphism}{{0BQN}{source}{}{sourcepointer-3}{}}
\newlabel{pione-proposition-universal-homeomorphism@cref}{{[originallemma][0][]0BQN}{[1][0][]source}{}{}{}}
\makeatother
\crefname{originaltheorem}{original theorem}{original theorems}
\Crefname{originaltheorem}{Original theorem}{Original theorems}
\makeatletter
\newlabel{theorem-topological-invariance}{{04DZ}{source}{}{sourcepointer-4}{}}
\newlabel{theorem-topological-invariance@cref}{{[originaltheorem][0][]04DZ}{[1][0][]source}{}{}{}}
\makeatother
\crefname{originallemma}{original lemma}{original lemmas}
\Crefname{originallemma}{Original lemma}{Original lemmas}
\makeatletter
\newlabel{pione-proposition-normal}{{0BQM}{source}{}{sourcepointer-5}{}}
\newlabel{pione-proposition-normal@cref}{{[originallemma][0][]0BQM}{[1][0][]source}{}{}{}}
\makeatother
\crefname{originallemma}{original lemma}{original lemmas}
\Crefname{originallemma}{Original lemma}{Original lemmas}
\makeatletter
\newlabel{pione-lemma-normal-pione-quotient-inertia}{{0BTD}{source}{}{sourcepointer-6}{}}
\newlabel{pione-lemma-normal-pione-quotient-inertia@cref}{{[originallemma][0][]0BTD}{[1][0][]source}{}{}{}}
\makeatother
\crefname{originalremark}{original remark}{original remarks}
\Crefname{originalremark}{Original remark}{Original remarks}
\makeatletter
\newlabel{section-galois-action-stalks}{{03QW}{source}{}{sourcepointer-7}{}}
\newlabel{section-galois-action-stalks@cref}{{[originalremark][0][]03QW}{[1][0][]source}{}{}{}}
\makeatother
\title{Finite type locally constant sheaves on an irreducible geometrically unibranch scheme are the finite representations of the etale fundamental group}
\author{The Stacks Project Authors}
\date{stacks-project:etale-cohomology.tex:local-snapshot-20261011; GNU FDL 1.2 or later}
\begin{document}
\maketitle
\noindent\textit{Editorial note.} All delivered mathematical text is the author's own, copied from the pinned Stacks Project chapter file etale-cohomology.tex; no statement or proof step was rewritten, reordered or removed.\par\medskip
\noindent\textit{Editorial note.} The proof is the author's complete argument. It opens by explaining why the companion lemma's proof does not apply here, and then supplies the three steps above inside the counted span.\par\medskip
\noindent\textit{Editorial note.} Every labelled reference occurring inside the counted spans is carried in assembly.external\_references as a non-load-bearing navigation pointer (same-chapter ones at their definition site, cross-chapter ones at the invocation inside the counted span), so no load-bearing mathematics is replaced by a stub. No bibliography citation occurs inside the counted spans.\par\medskip
\noindent\textit{Editorial note.} The checked chapter list for same-name collisions is the set of stacks chapters staged in this repository: chow, cohomology, divisors, etale-cohomology, flat, more-algebra, more-morphisms, morphisms, spaces-more-morphisms, varieties together with this card's chapter etale-cohomology.\par\medskip
\noindent\textit{Editorial note.} The preamble is the chapter's own native preamble excerpt already used for this source: the Xy-pic setup (native preamble lines 14-18) and the native theorem declarations with the mathematical macros (native preamble lines 154-217).\par\medskip
\noindent\textit{Allowed standard prerequisites (editorial).}\begin{itemize}
\item Geometrically unibranch scheme and its normalization. Same chapter together with the chapter Morphisms, the normalization morphism of a geometrically unibranch scheme, which is a universal homeomorphism. Supplies the reduction that replaces X by its normalization, using the invariance of the etale site under universal homeomorphisms.
\item Etale fundamental group and the equivalence for connected schemes. Same chapter, the companion result identifying finite locally constant sheaves of Lambda-modules on a connected scheme with finite Lambda-modules carrying a continuous action of pi\_1. Supplies the statement being adapted and the reason it cannot be quoted directly, namely that a finite Lambda-module need not have finite underlying set.
\item Finiteness of the global sections of a finite type locally constant sheaf on a connected normal scheme. Same chapter, the argument in the normal integral case identifying a finite type locally constant sheaf with a finite module over Lambda with continuous action of the fundamental group. Supplies the two directions constructed in the second and third paragraphs of the proof.
\end{itemize}\medskip
\ifdefined\hypertarget\hypertarget{sourcepointer-1}{}\fi
\noindent\textit{Source reference.} Original lemma 0DV5: sources/stacks-etale-cohomology/pinned-original-source.tex, source lines 10118--10118. Navigation pointer invoked at proof line 10222; the counted span contains the pointer only and the cited result is not reproduced inside it, so the pointer is registered as a non-load-bearing navigation stub in the established form for Stacks material. No bibliography citation is involved.\par\medskip
\ifdefined\hypertarget\hypertarget{sourcepointer-2}{}\fi
\noindent\textit{Source reference.} Original lemma 0GIQ: sources/stacks-etale-cohomology/pinned-original-source.tex, source lines 10228--10228. Navigation pointer invoked at proof line 10228; the counted span contains the pointer only and the cited result is not reproduced inside it, so the pointer is registered as a non-load-bearing navigation stub in the established form for Stacks material. No bibliography citation is involved.\par\medskip
\ifdefined\hypertarget\hypertarget{sourcepointer-3}{}\fi
\noindent\textit{Source reference.} Original lemma 0BQN: sources/stacks-etale-cohomology/pinned-original-source.tex, source lines 10230--10230. Navigation pointer invoked at proof line 10230; the counted span contains the pointer only and the cited result is not reproduced inside it, so the pointer is registered as a non-load-bearing navigation stub in the established form for Stacks material. No bibliography citation is involved.\par\medskip
\ifdefined\hypertarget\hypertarget{sourcepointer-4}{}\fi
\noindent\textit{Source reference.} Original theorem 04DZ: sources/stacks-etale-cohomology/pinned-original-source.tex, source lines 6354--6354. Navigation pointer invoked at proof line 10233; the counted span contains the pointer only and the cited result is not reproduced inside it, so the pointer is registered as a non-load-bearing navigation stub in the established form for Stacks material. No bibliography citation is involved.\par\medskip
\ifdefined\hypertarget\hypertarget{sourcepointer-5}{}\fi
\noindent\textit{Source reference.} Original lemma 0BQM: sources/stacks-etale-cohomology/pinned-original-source.tex, source lines 10241--10241. Navigation pointer invoked at proof line 10241; the counted span contains the pointer only and the cited result is not reproduced inside it, so the pointer is registered as a non-load-bearing navigation stub in the established form for Stacks material. No bibliography citation is involved.\par\medskip
\ifdefined\hypertarget\hypertarget{sourcepointer-6}{}\fi
\noindent\textit{Source reference.} Original lemma 0BTD: sources/stacks-etale-cohomology/pinned-original-source.tex, source lines 10249--10249. Navigation pointer invoked at proof line 10249; the counted span contains the pointer only and the cited result is not reproduced inside it, so the pointer is registered as a non-load-bearing navigation stub in the established form for Stacks material. No bibliography citation is involved.\par\medskip
\ifdefined\hypertarget\hypertarget{sourcepointer-7}{}\fi
\noindent\textit{Source reference.} Original remark 03QW: sources/stacks-etale-cohomology/pinned-original-source.tex, source lines 8516--8516. Navigation pointer invoked at proof line 10254; the counted span contains the pointer only and the cited result is not reproduced inside it, so the pointer is registered as a non-load-bearing navigation stub in the established form for Stacks material. No bibliography citation is involved.\par\medskip
\begin{lemma}
\label{lemma-locally-constant-on-connected-geom-unibranch}
Let $X$ be an irreducible, geometrically unibranch scheme.
Let $\overline{x}$ be a geometric point of $X$.
Let $\Lambda$ be a ring. There is an equivalence of categories
$$
\left\{
\begin{matrix}
\text{finite type, locally constant}\\
\text{sheaves of }\Lambda\text{-modules on }X_\etale
\end{matrix}
\right\}
\longleftrightarrow
\left\{
\begin{matrix}
\text{finite }\Lambda\text{-modules }M\text{ endowed}\\
\text{with a continuous }\pi_1(X, \overline{x})\text{-action}
\end{matrix}
\right\}
$$
\end{lemma}

\begin{proof}
The proof given in Lemma \ref{lemma-locally-constant-on-connected}
does not work as a finite $\Lambda$-module $M$ may not have a finite
underlying set.

\medskip\noindent
Let $\nu : X^\nu \to X$ be the normalization morphism. By Morphisms, Lemma
\ref{morphisms-lemma-normalization-geom-unibranch-univ-homeo}
this is a universal homeomorphism. By
Fundamental Groups, Proposition \ref{pione-proposition-universal-homeomorphism}
this induces an isomorphism
$\pi_1(X^\nu, \overline{x}) \to \pi_1(X, \overline{x})$
and by Theorem \ref{theorem-topological-invariance} we get
an equivalence of category between finite type, locally constant
$\Lambda$-modules on $X_\etale$ and on $X^\nu_\etale$.
This reduces us to the case where $X$ is an integral normal scheme.

\medskip\noindent
Assume $X$ is an integral normal scheme. Let $\eta \in X$ be the generic point.
Let $\overline{\eta}$ be a geometric point lying over $\eta$.
By Fundamental Groups, Proposition \ref{pione-proposition-normal}
have a continuous surjection
$$
\text{Gal}(\kappa(\eta)^{sep}/\kappa(\eta)) = \pi_1(\eta, \overline{\eta})
\longrightarrow
\pi_1(X, \overline{\eta})
$$
whose kernel is described in Fundamental Groups, Lemma
\ref{pione-lemma-normal-pione-quotient-inertia}.
Let $\mathcal{F}$ be a finite type, locally constant sheaf of $\Lambda$-modules
on $X_\etale$. Let $M = \mathcal{F}_{\overline{\eta}}$
be the stalk of $\mathcal{F}$ at $\overline{\eta}$.
We obtain a continuous action of $\text{Gal}(\kappa(\eta)^{sep}/\kappa(\eta))$
on $M$ by Section \ref{section-galois-action-stalks}.
Our goal is to show that this action factors through the
displayed surjection.
Since $\mathcal{F}$ is of finite type, $M$ is a finite $\Lambda$-module.
Since $\mathcal{F}$ is locally constant, for every $x \in X$
the restriction of $\mathcal{F}$ to $\Spec(\mathcal{O}_{X, x}^{sh})$
is constant. Hence the action of $\text{Gal}(K^{sep}/K_x^{sh})$
(with notation as in Fundamental Groups, Lemma
\ref{pione-lemma-normal-pione-quotient-inertia})
on $M$ is trivial. We conclude we have the factorization as desired.

\medskip\noindent
On the other hand, suppose we have a finite $\Lambda$-module $M$
with a continuous action of $\pi_1(X, \overline{\eta})$.
We are going to construct an $\mathcal{F}$ such that
$M \cong \mathcal{F}_{\overline{\eta}}$ as
$\Lambda[\pi_1(X, \overline{\eta})]$-modules.
Choose generators $m_1, \ldots, m_r \in M$.
Since the action of $\pi_1(X, \overline{\eta})$ on $M$ is
continuous, for each $i$ there exists an open subgroup
$N_i$ of the profinite group $\pi_1(X, \overline{\eta})$
such that every $\gamma \in H_i$ fixes $m_i$. We conclude that
every element of the open subgroup $H = \bigcap_{i = 1, \ldots, r} H_i$
fixes every element of $M$. After shrinking $H$ we may assume $H$
is an open normal subgroup of $\pi_1(X, \overline{\eta})$.
Set $G = \pi_1(X, \overline{\eta})/H$. Let $f : Y \to X$ be the
corresponding Galois finite \'etale $G$-cover.
We can view $f_*\underline{\mathbf{Z}}$ as a sheaf of
$\mathbf{Z}[G]$-modules on $X_\etale$. Then we just take
$$
\mathcal{F} =
f_*\underline{\mathbf{Z}} \otimes_{\underline{\mathbf{Z}[G]}} \underline{M}
$$
We leave it to the reader to compute
$\mathcal{F}_{\overline{\eta}}$.
We also omit the verification that this construction
is the inverse to the construction in the previous paragraph.
\end{proof}

\end{document}
